% ===========================================================================
%  LaTeX helper examples -- a cheat sheet, NOT part of the exposé.
%  Compiles on its own (open this file and save, or Ctrl+Alt+B).
%  Uses the same references.bib as expose.tex for the citation examples.
% ===========================================================================

\documentclass[11pt,a4paper]{article}

\usepackage[T1]{fontenc}
\usepackage{lmodern}
\usepackage[english]{babel}
\usepackage{microtype}
\usepackage{csquotes}
\usepackage[a4paper,margin=2.3cm]{geometry}
\usepackage{parskip}
\usepackage{amsmath,amssymb}
\usepackage{graphicx}
\usepackage{booktabs}
\usepackage[inline]{enumitem}
\usepackage{xcolor}
\usepackage[backend=biber,style=numeric-comp,sorting=none,
            maxbibnames=3,giveninits=true,doi=false]{biblatex}
\addbibresource{references.bib}
\usepackage[hidelinks]{hyperref}

\begin{document}

\section*{LaTeX Helper Examples}
\textcolor{red}{A cheat sheet for basic \LaTeX. Look at the source next to the
PDF to see how each part is written.}

\section{Text formatting}
\textbf{Bold}, \textit{italic}, \emph{emphasis}, \underline{underlined},
\texttt{monospace}, \textsc{Small Caps}, {\small small} and {\large large}.
Quotes: \enquote{double quotes}. Dashes: 1--2 (range), text --- aside.
A footnote.\footnote{Footnotes are numbered automatically.}
A link: \url{https://arxiv.org} or \href{https://arxiv.org}{named link}.
A forced line break\\ starts a new line; an empty line starts a new paragraph.

\section{Lists}
\begin{itemize}
  \item Bullet point
  \item Another one
  \begin{itemize}
    \item Nested bullet
  \end{itemize}
\end{itemize}

\begin{enumerate}
  \item Numbered item
  \item Second item
\end{enumerate}

\begin{description}
  \item[Term] Its explanation.
  \item[SAE] Sparse autoencoder.
\end{description}

Inline list: \begin{enumerate*}[label=(\roman*)] \item one, \item two,
\item three. \end{enumerate*}

\section{Mathematics}
Inline maths: the activation $x \in \mathbb{R}^d$ is reconstructed as
$\hat{x} = W_{\text{dec}} f(x) + b$.

Numbered equation, referenced later as Equation~\eqref{eq:sae}:
\begin{equation}
  \mathcal{L}(x) = \lVert x - \hat{x} \rVert_2^2 + \lambda \lVert f(x) \rVert_1
  \label{eq:sae}
\end{equation}

Unnumbered equation:
\[
  \cos(u, v) = \frac{u^\top v}{\lVert u \rVert \, \lVert v \rVert}
\]

Several aligned lines (\verb|&| marks the alignment point, \verb|\\| ends a line):
\begin{align}
  f(x)    &= \operatorname{ReLU}(W_{\text{enc}} x + b_{\text{enc}}) \\
  \hat{x} &= W_{\text{dec}} f(x) + b_{\text{dec}} \nonumber
\end{align}

\section{Tables}
Table~\ref{tab:example} uses \texttt{booktabs} rules.
\begin{table}[h]
  \centering
  \caption{An example table.}
  \label{tab:example}
  \begin{tabular}{l c r}
    \toprule
    Method  & Layers & AUROC \\
    \midrule
    Probe   & 12     & 0.91  \\
    SAE     & 12     & 0.87  \\
    \bottomrule
  \end{tabular}
\end{table}

\section{Figures}
Figure~\ref{fig:example} is a placeholder box. For a real image, put the file
next to this \texttt{.tex} file and use
\verb|\includegraphics[width=0.6\textwidth]{myplot.png}|.
\begin{figure}[h]
  \centering
  \fbox{\parbox[c][3cm][c]{0.5\textwidth}{\centering Your figure here}}
  \caption{An example figure.}
  \label{fig:example}
\end{figure}

\section{Citations and cross-references}
\begin{itemize}
  \item Numeric citation: \cite{bereska2024review}; several at once:
        \cite{engels2025notall,gurnee2025manifolds}.
  \item Author in the text: \textcite{bereska2024review} review the field.
  \item Cite with a page: \cite[p.~4]{bereska2024review}.
  \item Cross-reference: see Section~\ref{sec:helper-code} on
        page~\pageref{sec:helper-code}.
\end{itemize}
New sources go into \texttt{references.bib}; they appear in the reference list
only once they are cited.

\section{Code and comments}
\label{sec:helper-code}
Inline code: \verb|model.run_with_cache(tokens)|.
\begin{verbatim}
for layer in range(n_layers):
    acts = cache[f"blocks.{layer}.hook_resid_post"]
\end{verbatim}
% A line starting with % is a comment and does not appear in the PDF.

\printbibliography

\end{document}
